% The orbit-lattice paper: a margin in Weil's criterion, and the Riemann hypothesis as a lattice minimum. The claims
% table of Appendix A names the Lean theorem behind each numbered result, and lib/check/Check.lean prints their axioms.
% Remarks, conjectures and the numerical statements are not Lean theorems, and Section 5 lists them. The paper claims
% only its own steps: the results it uses from the literature are documented in formalisations/ and cited.
\documentclass[11pt]{article}
\usepackage[T1]{fontenc}
% pdfLaTeX needs the input encoding and the glyph declarations below; XeTeX reads Unicode directly.
\ifdefined\XeTeXversion\else\usepackage[utf8]{inputenc}\fi
\usepackage{amsmath,amssymb,amsthm}
\usepackage{booktabs}
\usepackage{longtable}
\usepackage{array}
\usepackage[hidelinks]{hyperref}
% The Unicode symbols of the Lean statements quoted in Section 5.
\ifdefined\XeTeXversion\else
\DeclareUnicodeCharacter{2203}{\ensuremath{\exists}}
\DeclareUnicodeCharacter{2200}{\ensuremath{\forall}}
\DeclareUnicodeCharacter{211D}{\ensuremath{\mathbb R}}
\DeclareUnicodeCharacter{2115}{\ensuremath{\mathbb N}}
\DeclareUnicodeCharacter{2124}{\ensuremath{\mathbb Z}}
\DeclareUnicodeCharacter{2194}{\ensuremath{\leftrightarrow}}
\DeclareUnicodeCharacter{2192}{\ensuremath{\to}}
\DeclareUnicodeCharacter{2208}{\ensuremath{\in}}
\DeclareUnicodeCharacter{2260}{\ensuremath{\ne}}
\DeclareUnicodeCharacter{2264}{\ensuremath{\le}}
\fi

\newtheorem{theorem}{Theorem}[section]
\newtheorem{proposition}[theorem]{Proposition}
\newtheorem{lemma}[theorem]{Lemma}
\newtheorem{corollary}[theorem]{Corollary}
\theoremstyle{definition}
\newtheorem{definition}[theorem]{Definition}
\newtheorem{conjecture}{Conjecture}
\newtheorem*{remark}{Remark}

\newcommand{\lean}[1]{\texttt{#1}}
\newcommand{\file}[1]{\nolinkurl{#1}}
\newcommand{\KH}{K}
\newcommand{\sinc}{\operatorname{sinc}}
\newcommand{\RH}{\mathrm{RH}}
\emergencystretch=3em

\title{A margin in Weil's criterion:\\ the Riemann hypothesis as a lattice minimum}
\author{Ritchie Comley}
\date{Draft, 13 November 2026}

\begin{document}
\maketitle

\begin{abstract}
Let $\KH(r) = \sum_\rho r^{\rho - 1/2}/(\rho(1 - \rho))$, the sum over the nontrivial zeros of the Riemann zeta
function, and let $G_N$ be the matrix $(\sqrt{mn}\,\KH(m/n))_{m, n \le N}$, whose entries can be computed from the
primes. By Weil's positivity criterion, in a discrete form of a theorem of Matsumoto and Suzuki, the Riemann hypothesis
holds if and only if every $G_N$ is positive semidefinite, and under it $G_N$ is the Gram matrix of the vectors $w_n =
(n^\rho/|\rho|)_\rho$. We prove two strengthenings, each equivalent to the Riemann hypothesis: every Gram--Schmidt pivot
of $w_1, w_2, \dots$ is at least the first, $\|w_1\|^2 = \KH(1) = 2 + \gamma - \log 4\pi$; and $w_1$ is a shortest
nonzero vector of the lattice spanned by $w_1, \dots, w_N$, for every $N$. The content is a margin: under the Riemann
hypothesis every pivot after the first is at least $0.0577 > 1.249\,\KH(1)$, uniformly, and $\pm w_1$ are the only
shortest vectors. The proof applies Weil's explicit formula to a band-limited test function at height of order $N$,
which makes the prime side exactly Landau's quadratic form; that form is bounded at the top vertex of the divisibility
graph by Schur's test with an explicit Mertens bound, and the archimedean term by an interval certificate. The proof is
formalised in Lean~4 with Mathlib and depends only on Lean's standard axioms. Neither statement survives its natural
analogues: both fail for every curve over a finite field, through an integral null vector that the Frobenius forces,
and both fail for a positive proportion of real Dirichlet characters, at the first step. The margin is a property of the
constants of $\zeta$, not of the structure that gives the Riemann hypothesis elsewhere.
\end{abstract}

% TODO (draft): expand the proof sketches of Section 4 into full proofs; decide whether Section 8 stays in this paper;
% read Bombieri (2003), Broughan's book and Mazhouda (2013) before the novelty claims of Section 7 are final.

\section{Introduction}\label{sec:intro}

Let $\rho$ run over the nontrivial zeros of the Riemann zeta function, counted with multiplicity. Weil's positivity
criterion \cite{weil52} says that the Riemann hypothesis (RH) holds if and only if a quadratic functional, built from
the primes and the archimedean place through the explicit formula, is nonnegative on a suitable space of test
functions. This paper restricts it to one orbit. For $r > 0$ let
\begin{equation}\label{eq:K}
\KH(r) = \sum_\rho \frac{r^{\rho - 1/2}}{\rho(1 - \rho)}, \qquad G(m, n) = \sqrt{mn}\,\KH(m/n),
\end{equation}
and let $G_N = (G(m, n))_{m, n \le N}$ and $Q_N(x) = \sum_{m, n \le N} x_m x_n G(m, n)$ for real $x$. The series
converges absolutely, $\KH(1/r) = \KH(r)$, and
\begin{equation}\label{eq:d1}
d_1 := \KH(1) = \sum_\rho \frac{1}{\rho(1 - \rho)} = 2 + \gamma - \log 4\pi = 0.04619\ldots,
\end{equation}
$\gamma$ being Euler's constant. Through the explicit formula $\KH$ is a function of Chebyshev's $\psi$
(Proposition~\ref{prop:kernel}), so $G_N$ can be computed from the primes without the zeros. Under RH, with $\rho =
\tfrac12 + i\gamma_\rho$, $G(m, n) = \sum_\rho m^\rho\,\overline{n^\rho}/|\rho|^2$: $G_N$ is the Gram matrix of the
vectors $w_n = (n^\rho/|\rho|)_\rho$ in $\ell^2$ of the zeros, and $Q_N(x) = \sum_\rho |\sum_{n \le N} x_n n^\rho|^2 /
|\rho|^2$. Conversely, RH holds if and only if every $Q_N$ is nonnegative at the integer vectors
(Proposition~\ref{prop:positivity}). That positivity layer is known in substance: $\KH$ is the function $H_1$ of
Matsumoto and Suzuki, who prove that RH holds exactly when $H_1(e^t) - H_1(1)$ is a screw function on $\mathbb R$
\cite[Theorem~1.3]{matsumoto-suzuki}, and the $2 \times 2$ minors of $G_N$ are Nicolas's bound $|W(x)| \le \beta$, with
$W = \KH$ under RH and $\beta = d_1$ \cite{nicolas}.

We ask for more than positivity. Let $d_k$ be the $k$-th Gram--Schmidt pivot of the orbit, the least value of
$Q_k(\delta_k - \sum_{n < k} c_n \delta_n)$ over real $c$: under RH, the squared distance of $w_k$ from the span of
$w_1, \dots, w_{k - 1}$. Consider
\begin{itemize}
\item[(A)] $d_k \ge d_1$ for every $k \ge 1$ (the pivot form);
\item[(B)] $Q_N(a) \ge d_1$ for every $N$ and every nonzero $a \in \mathbb Z^N$ (the shortest-vector form): under RH,
$w_1$ is a shortest nonzero vector of the lattice spanned by $w_1, \dots, w_N$.
\end{itemize}

\begin{theorem}\label{thm:main}
The following are equivalent: \textup{(i)} the Riemann hypothesis; \textup{(ii)} \textup{(A)}; \textup{(iii)}
\textup{(B)}.
\end{theorem}

The implications from (A) and from (B) to RH are formal: each makes every $Q_N$ nonnegative at the integer vectors. The
content of Theorem~\ref{thm:main} is the converse, a margin that holds uniformly:

\begin{theorem}\label{thm:margin}
Assume RH. For every $N \ge 2$ and every real $x$,
\[
Q_N(x) \ge c_0\, x_N^2, \qquad c_0 = \frac{999}{1000} \cdot \frac{2}{45} \cdot \frac{13}{10} = 0.05772.
\]
Hence $d_k \ge c_0 > 1.249\,d_1$ for every $k \ge 2$; every integer vector with a nonzero coordinate at some $n \ge 2$
has $Q_N(a) \ge c_0$; and $\pm\delta_1$ are the only nonzero integer vectors with $Q_N(a) \le d_1$.
\end{theorem}

The true values are larger. Numerically $d_2/d_1 = 1.7184\ldots$, certified in ball arithmetic, and the second minimum of
the lattice, the least $Q_N(a)/d_1$ over integer $a \ne 0, \pm\delta_1$, is $1.9387\ldots$ for $N \le 30$, at $a =
\delta_1 + \delta_2$ (Section~\ref{sec:numerics}). Under RH the pivots grow like $\log k$
(Proposition~\ref{prop:growth}). In the flat model, where the zeros are replaced by their density, the Gram matrix is
$d_1 \min(m, n)$ and every pivot equals $d_1$ (Proposition~\ref{prop:flat}); the zeros of $\zeta$ push every pivot
after the first above $d_1$, and keep the lattice's shortest vector unique.

\subsection*{What is new}

As far as a literature search reaches \cite{repo-prior-art}, neither (A) nor (B) has been stated, no source measures the
Gram--Schmidt pivots of an explicit-formula Gram matrix against the first pivot, and none states RH through a lattice
minimum. The nearest results are Lagarias's Gram matrix of Weil's scalar product on Li's functions
\cite[Theorem~3.1]{lagarias07}, the same pairing on another orbit of $1/s$, whose pivots collapse to $0$ under RH
(Section~\ref{sec:li}); Yoshida's and Bombieri's non-degeneracy and minimisation results for Weil's form on windows
\cite{yoshida92, bombieri00, bombieri03}, which bound squared norms rather than pivots; and Hallouin and Perret's Gram
matrix of the Frobenius orbit of a curve \cite{hallouin-perret}, where the margin fails. Section~\ref{sec:related}
places the paper in that literature. We claim the margin, Theorem~\ref{thm:margin}, and the equivalences it gives; the
positivity layer, the kernel and the constant $d_1$ are not ours.

\subsection*{The analogues}

The same construction can be made for a curve over a finite field, where RH is a theorem, and for a Dirichlet
$L$-function, where GRH is expected. Both forms fail there (Section~\ref{sec:analogues}). For a curve the zeros repeat
vertically, and the Frobenius forces a null vector whose entries are the coefficients of its minimal polynomial: an
integer vector with $Q = 0$. For characters the first step decides: every orbit Gram matrix has $G_{22} = 2G_{11}$, so on
$w_1, w_2$ the two forms coincide, and a positive proportion of real characters fail at $k = 2$. The forms part only on
the geometric orbit of the Frobenius graphs, whose integer span is the lattice $\mathbb Z[z]/(m)$: there (A) fails for
every curve and (B) holds, by integrality, whenever $m$ is irreducible. For $\zeta$ no integral structure is known, and
(B) rests on the analytic margin alone. So the margin is a property of the constants of $\zeta$, its conductor $1$,
$d_1 = 0.046$ and the absence of zeros below height $14.13$, and Theorem~\ref{thm:main} is a form of RH, not a route to
it.

\subsection*{The proof and its formalisation}

The proof of Theorem~\ref{thm:margin} (Section~\ref{sec:proof}) applies Weil's explicit formula to $h(t) = |D(t)|^2
\varphi(t/N)/N^2$, with $D(t) = \sum_{n \le N} x_n \sqrt n\, n^{-it}$ and a fixed test function $\varphi$ whose Fourier
transform is supported in $[-\tfrac45, \tfrac45]$. Distinct frequencies $\log(jn)$ and $\log j'$ are at least $4/(5N)$
apart, so the prime side is exactly Landau's quadratic form. A minorant of the archimedean term, band-limited to $[-1,
1]$, keeps only the diagonal, and an interval certificate checks it against the digamma function. The last step is
arithmetic: under the diagonal $\log N + d$, Landau's form leaves the top vertex $N$ of the divisibility graph
$\tfrac{13}{10}$ of its weight, by Schur's test with an explicit bound for Mertens' sums
(Proposition~\ref{prop:top}). The whole proof is formalised in Lean~4 \cite{lean4} with Mathlib \cite{mathlib}, in the
repository of the paper \cite{repo}, and depends only on the standard axioms (Section~\ref{sec:lean}).
Appendix~\ref{app:claims} names the Lean theorem behind each numbered result. Remarks, conjectures and numerical
statements are not Lean theorems.

\subsection*{Organisation}

Section~\ref{sec:setup} sets up the kernel, the orbit and the two forms. Section~\ref{sec:results} states the main
results and proves the easy implications, and Section~\ref{sec:proof} proves the margin. Section~\ref{sec:lean}
describes the formalisation and Section~\ref{sec:analogues} the analogues. Section~\ref{sec:related} discusses the
literature, and Section~\ref{sec:lambda} the smallest eigenvalue of $G_N$ and some questions.

\section{The kernel, the orbit and the two forms}\label{sec:setup}

\begin{proposition}\label{prop:kernel}
The series \eqref{eq:K} converges absolutely for every $r > 0$, $\KH(1/r) = \KH(r)$, $\KH(1) = 2 + \gamma - \log 4\pi$,
and for $x \ge 1$, with $u = \log x$,
\begin{multline*}
\KH(x) = \cosh(\tfrac u2)(2 + \gamma - \log 4\pi) + \sinh(\tfrac u2)(\log \pi - 2u) + A(u) \\
+ \sqrt x \sum_{n \le x} \frac{\Lambda(n)}{n} - \frac{\psi(x)}{\sqrt x},
\end{multline*}
where $\psi(x) = \sum_{n \le x} \Lambda(n)$ and $A(u) = \frac{1}{2\pi} \int_{\mathbb R} H_u(r) \operatorname{Re}
\psi_\Gamma(\tfrac14 + \tfrac{ir}2)\,dr$, with $H_u(t) = (\cos(tu) - \cosh(u/2))/(t^2 + \frac14)$ and $\psi_\Gamma$ the
digamma function.
\end{proposition}

The value $\KH(1) = \sum_\rho 1/(\rho(1 - \rho))$ comes from the Hadamard product (see \cite[Chapter~12]{davenport}); the
equivalence of RH with $\sum_\rho 1/|\rho|^2 = 2 + \gamma - \log 4\pi$ is Burnol's remark \cite{burnol98}, restated by
Gun, Murty and Rath \cite{gun-murty-rath}. The formula for $\KH(x)$ is Weil's explicit formula
(Lemma~\ref{lem:explicit}) for the test function $H_u$, whose poles at $\pm i/2$ are removable
(\file{formalisations/ball-rate.md}).

\begin{proposition}\label{prop:positivity}
The Riemann hypothesis holds if and only if $Q_N(a) \ge 0$ for every $N$ and every $a \in \mathbb Z^N$. Under RH there
are vectors $w_n$, $n \ge 1$, in a complex Hilbert space with $\langle w_m, w_n\rangle = G(m, n)$, so that $Q_N(x) =
\|\sum_{n \le N} x_n w_n\|^2$.
\end{proposition}

The vectors $w_n = (n^\rho/|\rho|)_\rho$ in $\ell^2$ of the zeros realise $G$; in the library a realisation is a family
indexed by the positive rationals on which multiplication acts by similitudes. The converse uses the $2 \times 2$ minors
of $G$, at the squares $(p/q)^2$, which make $2\KH(1) + \KH(r) + \KH(1/r)$ nonnegative on a dense set of $r$, and
Landau's theorem for the Laplace transform of a nonnegative function. As said in the
introduction, this layer is a discrete form of Matsumoto and Suzuki's theorem \cite{matsumoto-suzuki}: $\KH$ is their
$H_1$, and the kernel $\KH(e^{t - u}) - \KH(e^t) - \KH(e^{-u}) + \KH(1)$ of their screw function is the Gram matrix of
the differences $w_n/\sqrt n - w_1$.

\begin{definition}\label{def:forms}
For $k \ge 1$ let $d_k = \inf_c Q_k(\delta_k - \sum_{n < k} c_n \delta_n)$, the infimum over real $c_1, \dots, c_{k -
1}$; $d_1 = \KH(1)$. The \emph{pivot form} (A) is the statement $d_k \ge d_1$ for every $k \ge 1$. The
\emph{shortest-vector form} (B) is the statement $Q_N(a) \ge d_1$ for every $N \ge 1$ and every $a \in \mathbb Z^N
\setminus \{0\}$.
\end{definition}

(B) asks less than (A). If $a_k$ is the last nonzero coordinate of an integer vector $a$, then $a_k^2 \ge 1$ and $Q(a)
\ge a_k^2 d_k \ge d_k$, by homogeneity. The flat model is the equality case of both:

\begin{proposition}\label{prop:flat}
Let $\KH_0(x) = e^{-|\log x|/2} = \frac{1}{2\pi} \int_{\mathbb R} \cos(r \log x)/(r^2 + \frac14)\,dr$. Then
$\sqrt{ab}\,\KH_0(a/b) = \min(a, b)$ for $a, b > 0$, and vectors with Gram matrix $c \min(m, n)$ have every Gram--Schmidt
pivot equal to $c$.
\end{proposition}

So replacing the zeros by Lebesgue measure with the weight $1/(\frac14 + t^2)$ gives the covariance of Brownian motion
at integer times, whose increments are orthogonal of equal norm. In that model (A) and (B) hold with equality, and the
increments $\delta_n - \delta_{n - 1}$ are all shortest vectors. The orthogonal increments are Lubinsky's orthogonal
Dirichlet polynomials for the arctangent density \cite{lubinsky14}, in Dimitrov and Oliveira's form
\cite[\S2.3]{dimitrov-oliveira}.

\section{The main results}\label{sec:results}

\begin{theorem}[Theorem~\ref{thm:main}]\label{thm:main2}
RH holds if and only if \textup{(A)} holds, and if and only if \textup{(B)} holds.
\end{theorem}

\begin{proof}
If (A) holds, then every pivot is at least $d_1 > 0$, so by induction on $N$, $Q_N(x) \ge d_N x_N^2 + Q_{N -
1}(x') \ge 0$ for every real $x$, where $x'$ is $x$ without its last coordinate; RH follows by
Proposition~\ref{prop:positivity}. If (B) holds, then $Q_N(a) \ge 0$ at every integer vector, and RH follows by the
same proposition. Conversely, assume RH. For $k = 1$, $Q_1(\delta_1) = d_1$. For $k \ge 2$, Theorem~\ref{thm:margin} gives
$Q_k(\delta_k - \sum_{n < k} c_n \delta_n) \ge c_0 > d_1$, since $d_1 < 0.0539$. So (A) holds. If $a \in \mathbb
Z^N$ is nonzero, either $a$ is supported at $1$, and $Q_N(a) = a_1^2 d_1 \ge d_1$, or some $a_n$, $n \ge 2$, is
nonzero, and $Q_N(a) \ge c_0 > d_1$ by Theorem~\ref{thm:margin}. So (B) holds.
\end{proof}

\begin{corollary}\label{cor:lattice}
Assume RH. For every $N \ge 1$, the nonzero integer vectors $a$ with $Q_N(a) \le d_1$ are $\pm\delta_1$, and every other
nonzero integer vector has $Q_N(a) \ge c_0 = 1.2496\ldots d_1$. In every realisation, the lattice spanned by $w_1, \dots,
w_N$ has minimum $\|w_1\|^2 = d_1$, attained at $\pm w_1$ only, and every other lattice vector is longer than $w_1$ by
a factor at least $\sqrt{1.249} = 1.117\ldots$.
\end{corollary}

\begin{proposition}\label{prop:growth}
Assume RH. For every $N \ge 1$ and every real $x$, $Q_N(x) \ge (\tfrac12 \log N + \log 4 + 2)\, x_N^2 / (22 e^8)$. So
$d_N \to \infty$.
\end{proposition}

Proposition~\ref{prop:growth} comes from an earlier sampling theorem of the library (\file{formalisations/ball-rate.md},
the sections on \lean{BallSharp} and \lean{BallPivot}); its constant passes $d_1$ only once $\log N \ge 7100$, and
Theorem~\ref{thm:margin} replaces it at every $N$. The upper bound $d_N \le \log(N/2\pi) + 1 - \gamma + o(1)$ shows that
$\log N$ is the true order; it is not formalised.

\subsection{Numerical values}\label{sec:numerics}

\begin{remark}
These are computations, not Lean theorems. In ball arithmetic $d_2/d_1 = 1.71841346690985\ldots$
(\texttt{scripts/ball\_tests/hunt/k2.py}), and $d_k > d_1$ for $2 \le k \le 10^4$ is certified with every rounding
bounded, without RH (\texttt{scripts/ball\_certificate.py}). An independent computation in ball arithmetic up to $k =
500$ puts the minimum of $d_k/d_1$ at $k = 2$, with $d_{10}/d_1 = 6.35$, $d_{100}/d_1 = 35.67$ and $d_{500}/d_1 = 74.93$
(\texttt{scripts/ball\_tests/hunt/cert\_zeta.py}). A Fincke--Pohst enumeration
of the integer vectors with $Q_N(a) \le 3d_1$ finds, for $N = 12$ and $N = 30$, that the second minimum of the lattice
is $Q_N(\delta_1 + \delta_2)/d_1 = 1.93870544505279\ldots$, the value certified in ball arithmetic
(\texttt{scripts/ball\_tests/shortest/shortest.py}).
\end{remark}

\section{Proof of the margin}\label{sec:proof}

Throughout, $N \ge 2$, $x$ is real, $c_n = x_n \sqrt n$, $D(t) = \sum_{n \le N} c_n n^{-it}$, and $\Lambda$ is von
Mangoldt's function. For a test function $h$ let $g_h(u) = \frac{1}{2\pi} \int_{\mathbb R} h(r) \cos(ru)\,dr$.

\begin{lemma}[Weil's explicit formula]\label{lem:explicit}
Let $h$ be even, holomorphic on $|\operatorname{Im} t| \le 1$ with $|h(t)| \le C/(1 + (\operatorname{Re} t)^2)$ there,
and real on $\mathbb R$. Then, with $t_\rho = (\rho - \frac12)/i$,
\begin{multline*}
\sum_\rho h(t_\rho) = h(\tfrac i2) + h(-\tfrac i2) - g_h(0) \log \pi \\
+ \frac{1}{2\pi} \int_{\mathbb R} h(r) \operatorname{Re} \psi_\Gamma(\tfrac14 + \tfrac{ir}{2})\,dr - 2 \sum_{n \ge 1}
\frac{\Lambda(n)}{\sqrt n}\, g_h(\log n).
\end{multline*}
\end{lemma}

This is the explicit formula of Guinand and Weil \cite{weil52}; the library proves it from the Hadamard product of $\xi$
(\file{formalisations/ball-rate.md}, \file{formalisations/weil-positivity.md}).

\subsection{The test function}

Let $\sinc x = (\sin x)/x$, $\mathcal K(x) = \sinc^2 x\,(1 + \frac13 \sin^2 x)$, and
\[
\varphi(s) = \frac{999}{1000} \cdot \frac{1 - \mathcal K(s/5)}{s^2}, \qquad \hat\varphi(v) = \int_{\mathbb R}
\varphi(s) \cos(sv)\,ds.
\]

\begin{lemma}\label{lem:phi}
\textup{(i)} $\hat\varphi(v) = \frac{999}{1000}\pi\big(\frac{25}{9}((\frac25 - |v|)^+)^3 - \frac{25}{144}((\frac45 -
|v|)^+)^3\big)$; in particular $\hat\varphi$ vanishes for $|v| \ge \frac45$ and $\hat\varphi(0) = \frac{999}{1000}
\cdot \frac{4\pi}{45}$. \textup{(ii)} $0 \le \varphi(s)$ and $\varphi(s)(s^2 + \frac1{16}) \le 1$ for real $s$.
\textup{(iii)} $\varphi$ extends to an entire function $\varphi_{\mathbb C}$ with $|\varphi_{\mathbb C}(z)| \le C/(1 +
(\operatorname{Re} z)^2)$ on $|\operatorname{Im} z| \le 1$, and $|\varphi_{\mathbb C}(iy)| \le y^2/3000$ for $0 < y
\le \frac14$.
\end{lemma}

\begin{proof}[Sketch]
$\varphi$ is a combination of the atoms $(1 - \sinc^2(\beta s))/s^2$ at $\beta = \frac15$ and $\frac25$, whose transforms
are cubic splines supported in $|v| \le 2\beta$. Part (ii) follows from $0 \le \mathcal K \le 1$ and an elementary
bound for $1 - \mathcal K$; in (iii) the terms of order $y^2$ of $1 - \mathcal K(iy/5)$ cancel.
\end{proof}

Let $\Phi_N(t) = \varphi(t/N)/N^2$ and $h = |D|^2 \Phi_N$, so that $h(t) = D(t) D(-t) \Phi_N(t)$ extends to an even
entire function to which Lemma~\ref{lem:explicit} applies. By Lemma~\ref{lem:phi}(ii), $\Phi_N(t) \le 1/(t^2 + N^2/16)
\le 1/(\frac14 + t^2)$ for real $t$ and $N \ge 2$, so under RH, where $t_\rho = \gamma_\rho$ is real and $D(\gamma_\rho)
= \sum_n x_n n^{1 - \rho}$,
\begin{equation}\label{eq:zero-side}
\sum_\rho h(t_\rho) = \sum_\rho \Big|\sum_{n \le N} x_n n^\rho\Big|^2 \Phi_N(\gamma_\rho) \le Q_N(x).
\end{equation}

\subsection{The prime side}

\begin{lemma}\label{lem:prime}
$2 \sum_n \Lambda(n) n^{-1/2} g_h(\log n) = \frac{\hat\varphi(0)}{2\pi N} L(c)$, where $L(c) = 2 \sum_{n \le N}
\Lambda(n) n^{-1/2} e_n(c)$ is Landau's form and $e_n(c) = \sum_{jn \le N} c_j c_{jn}$.
\end{lemma}

\begin{proof}
Expanding $|D(r)|^2 = \sum_{j, j'} c_j c_{j'} \cos(r \log(j/j'))$ and using $\int \Phi_N(r) e^{irv}\,dr =
\hat\varphi(Nv)/N$, $g_h(\log n)$ is a combination of the values $\hat\varphi(N(\log(jn) - \log j'))$. If $jn \ne j'$
with $j, j' \le N$, then $|\log(jn) - \log j'| \ge \log(1 + \frac1N) \ge \frac{2}{2N + 1} \ge \frac{4}{5N}$, by
$\log\frac{1 + u}{1 - u} \ge 2u$, and $\hat\varphi$ vanishes there. Only the exact matches $jn = j'$ remain, and
$g_h(\log n) = \hat\varphi(0) e_n(c)/(2\pi N)$.
\end{proof}

\subsection{The archimedean term}

Let $\lambda(s) = \frac{23}{100} + \frac{693}{1000} \sum_{k < 30} (1 - \sinc^2(s/(10 \cdot 2^k)))$, a staircase whose
steps of height $0.693 < \log 2$ at the scales $10 \cdot 2^k$ follow $\log s$, and let $r(s) = \sum_{j = 1}^{9} b_j
(\sinc^2(\frac{s - \gamma_j}{2}) + \sinc^2(\frac{s + \gamma_j}{2}))$, nine Fejér bumps with explicit rational $b_j > 0$
and $0 < \gamma_j \le \frac{17}{2}$, $\sum_j b_j = 0.004561\ldots$, which pay for the heights $s < 12$. Let $\mu =
\varphi\lambda - r$.

\begin{lemma}\label{lem:minorant}
\textup{(i)} $\int \mu(s) \cos(sv)\,ds = 0$ for $|v| \ge 1$, and $\int \mu \ge \frac78 \hat\varphi(0)$.
\textup{(ii)} For every $N \ge 2$ and every real $s$, $\mu(s) \le \varphi(s)\big(\operatorname{Re}
\psi_\Gamma(\frac14 + \frac{iNs}{2}) - \log \pi - \log N\big)$.
\end{lemma}

\begin{proof}[Sketch]
(i) A step at scale $A \ge 10$ widens the support of $\hat\varphi$ by at most $2/A \le \frac15$, and each bump has
transform supported in $[-1, 1]$. The mass is exact: by $\sinc^2 y = \frac14 \int (2 - |w|)^+ \cos(wy)\,dw$ and Fubini,
$\int \mu = \hat\varphi(0)(\frac{23}{100} + \frac{693}{1000} \sum_k \kappa(10 \cdot 2^k)) - 4\pi \sum_j b_j$ with
$\kappa(A) = \frac{15}{2A} - \frac{175}{8A^2} + \frac{375}{16A^3}$, which is $0.8804\,\hat\varphi(0)$.
(ii) Stirling's formula with two steps of the recurrence gives $\operatorname{Re} \psi_\Gamma(\frac14 + \frac{iNs}2) -
\log \pi - \log N \ge g(s) = \log(s/2\pi) - \bar e(2s)$ for every $N \ge 2$ at once, with a rational function $\bar e$
that decreases, so it suffices that $\varphi(\lambda - g) \le r$. That is an interval certificate in exact rational
arithmetic: $\lambda$ is bounded above through $1 - \sinc^2 x \le \min(1, \frac{x^2}3 - \frac{2x^4}{45} +
\frac{x^6}{315})$, $\varphi$ above and $r$ below through $\sinc x \ge 1 - \frac{x^2}6 + \frac{x^4}{120} -
\frac{x^6}{5040}$, and $\log y$ below by $2v + \frac23 v^3$, $v = \frac{y - 1}{y + 1}$. Above $s = 20$, writing $s = 10
\cdot 2^j y$ with $y \in [1, 2)$ reduces every $j$ to the same $400$ cells of $y$; $[12, 20)$ and $[\frac12, 12)$ take
$320$ and $460$ cells, and $(0, \frac12]$ follows from $\varphi(s) \le s^2/3000$. The least margin is $0.0042$ above $s =
12$ and $2.4\%$ below.
\end{proof}

\begin{lemma}\label{lem:pole}
$h(\frac i2) + h(-\frac i2) \ge -\sum_{n \le N} c_n^2/(6000N)$.
\end{lemma}

\begin{proof}[Sketch]
$h(\pm\frac i2) = D(\frac i2) D(-\frac i2) \varphi_{\mathbb C}(\frac{i}{2N})/N^2$; bound $|D(\pm\frac i2)|$ by
Cauchy--Schwarz and $\varphi_{\mathbb C}$ by Lemma~\ref{lem:phi}(iii) at $y = 1/(2N) \le \frac14$.
\end{proof}

\begin{proposition}\label{prop:form}
Assume RH. For $N \ge 2$ and real $x$,
\[
Q_N(x) \ge \frac{\hat\varphi(0)}{2\pi N}\Big(\big(\log N + \tfrac78\big) \sum_{n \le N} c_n^2 - L(c)\Big) - \frac{1}{6000N}
\sum_{n \le N} c_n^2 .
\]
\end{proposition}

\begin{proof}
By \eqref{eq:zero-side} and Lemma~\ref{lem:explicit}, $Q_N(x)$ is at least the pole terms, plus the archimedean term
$\frac{1}{2\pi} \int h(r)(\operatorname{Re} \psi_\Gamma(\frac14 + \frac{ir}2) - \log \pi)\,dr$, minus the prime side.
Substituting $r = Ns$, the archimedean term is $\frac{1}{2\pi N} \int |D(Ns)|^2 \varphi(s)(\log N + (\operatorname{Re}
\psi_\Gamma(\frac14 + \frac{iNs}2) - \log \pi - \log N))\,ds$, which by Lemma~\ref{lem:minorant}(ii) is at least
$\frac{1}{2\pi N} \int |D(Ns)|^2 (\varphi(s) \log N + \mu(s))\,ds$. Distinct frequencies $\log j \ne \log j'$, $j, j' \le
N$, are at least $1/N$ apart, so against $|D(N\cdot)|^2$ the transforms of $\varphi$ and $\mu$, supported in $[-\frac45,
\frac45]$ and $[-1, 1]$, keep only the diagonal: the term is $\frac{1}{2\pi N}(\hat\varphi(0) \log N + \int \mu) \sum
c_n^2 \ge \frac{\hat\varphi(0)}{2\pi N}(\log N + \frac78) \sum c_n^2$. Lemmas~\ref{lem:prime} and~\ref{lem:pole} give
the prime side and the poles.
\end{proof}

\subsection{The top vertex}

Let $M(y) = \sum_{n \le y} \Lambda(n)/n$.

\begin{lemma}\label{lem:mertens}
$M(y) \le \log y + \frac{81}{100}$ for every integer $y \ge 1$. Moreover $M(y) \le \log y + (\log 3 - 1) + (1 + \log
y)/y$ for $1 \le y \le 2000$.
\end{lemma}

\begin{proof}[Sketch]
$\psi(n) = \log \operatorname{lcm}(1, \dots, n)$, and $\operatorname{lcm}(1, \dots, n) \le 3^n$ for $n \le 2000$, by one
kernel computation. With $yM(y) \le \sum_{m \le y} \log m + \psi(y)$ and Stirling's bound this gives the second
inequality, and Chebyshev's $\psi(y) \le y \log 4 + 2\sqrt y \log y$ gives the first beyond $y = 1600$.
\end{proof}

\begin{proposition}\label{prop:top}
For every $N \ge 2$, every $d \ge \frac{17}{20}$ and every real $c$,
\[
\big(\log N + d\big) \sum_{n \le N} c_n^2 - L(c) \ge \tfrac{13}{10}\, c_N^2 .
\]
\end{proposition}

\begin{proof}[Sketch]
Only prime powers $q$ contribute to $L(c)$, through the edges $(j, jq)$ of the divisibility graph on $[1, N]$. Split off
the edges into the top vertex: $L(c) = L(c') + 2c_N \sum_{q \mid N} \Lambda(q) q^{-1/2} c_{N/q}$, where $c'$ is $c$ with
$c_N = 0$. Schur's test, row by row, bounds $L(c')$ by $\sum_{j < N} c_j^2 (M(\lfloor (N - 1)/j \rfloor) + \log j)$, so
row $j$ keeps the part $\rho_j = \log N + d - M(\lfloor (N - 1)/j \rfloor) - \log j \ge 0$ of its diagonal. At a
parent $j = N/q$ of the top vertex, $\lfloor (N - 1)/j \rfloor = q - 1$ and $\rho_j = \operatorname{room}_d(q) = d +
\log q - M(q - 1)$. Completing the square at each parent costs $c_N^2 \Lambda(q)^2/(q \operatorname{room}_d(q))$, and the
parents are distinct, so by $\sum_{q \mid N} \Lambda(q) = \log N$ the coefficient of $c_N^2$ is at least
\[
d + \sum_{q \mid N} \Lambda(q)\Big(1 - \frac{\Lambda(q)}{q \operatorname{room}_d(q)}\Big).
\]
Every term is nonnegative, and a prime $p \mid N$ contributes at least $\frac{45}{100}$: $\operatorname{room}_d(q)$ is
computed exactly for $q = 2, 3, 4, 5, 7$, it is at least $d - 0.52$ on $[8, 2001]$ and $d - 0.43$ on $[11, 2001]$ by the
second bound of Lemma~\ref{lem:mertens}, and at least $d - 0.81$ everywhere by the first. So the coefficient is at least
$\frac{17}{20} + \frac{45}{100} = \frac{13}{10}$.
\end{proof}

\subsection{Proof of Theorem~\ref{thm:margin}}

Since $\hat\varphi(0)/(2\pi) = \frac{999}{1000} \cdot \frac{2}{45}$, the pole term of Proposition~\ref{prop:form} is
exactly $\frac{\hat\varphi(0)}{2\pi N} \cdot \frac{5}{1332} \sum c_n^2$. So under RH
\[
Q_N(x) \ge \frac{\hat\varphi(0)}{2\pi N}\Big(\big(\log N + d\big) \sum_{n \le N} c_n^2 - L(c)\Big), \qquad d = \frac78
- \frac{5}{1332} \ge \frac{17}{20},
\]
and Proposition~\ref{prop:top} gives $Q_N(x) \ge \frac{\hat\varphi(0)}{2\pi N} \cdot \frac{13}{10} c_N^2 =
\frac{999}{1000} \cdot \frac{2}{45} \cdot \frac{13}{10}\, x_N^2$, as $c_N^2 = N x_N^2$. If $a$ is an integer vector
whose last nonzero coordinate is $a_k$, $k \ge 2$, then $Q_N(a) = Q_k(a) \ge c_0 a_k^2 \ge c_0$. Finally $d_1 < 0.0539
< c_0$, and $c_0/d_1 = 1.2496\ldots$ by \eqref{eq:d1}. \qed

\begin{remark}
A numerical design with a test function whose transform is supported in $[-\frac23, \frac23]$, steps that follow
$\log s$ more closely and a bound for the top vertex by $\Lambda(N)$ at each $N$ gives $d_k \ge 1.449\,d_1$ for every $k
\ge 2$ (\texttt{scripts/pivot\_push.py}); it is not formalised. The binding case of both is $k = 2$, where the true value
is $1.718\,d_1$.
\end{remark}

\section{The formalisation}\label{sec:lean}

Appendix~\ref{app:claims} names the Lean theorems behind each numbered result, in \file{lib/} of the repository
\cite{repo}; \file{formalisations/ball-rate.md} documents them, section by section. The Lean statements of
Theorem~\ref{thm:main} are
{\footnotesize
\begin{verbatim}
theorem rh_iff_orbitLattice :
    RiemannHypothesis ↔ ∀ (k : ℕ) (c : ℕ → ℝ), 1 ≤ k →
      BallWeil.gram 1 1 ≤
        BallWeil.form k (fun n => if n = k then 1 else -c n)

theorem rh_iff_shortest :
    RiemannHypothesis ↔ ∀ (N : ℕ) (a : ℕ → ℤ),
      (∃ n ∈ Icc 1 N, a n ≠ 0) →
        BallWeil.gram 1 1 ≤ BallWeil.form N (fun n => (a n : ℝ))
\end{verbatim}
}
\noindent in the namespace \lean{PivotSharp}, where \lean{RiemannHypothesis} is Mathlib's statement and
\lean{BallWeil.gram m n} is $\sqrt{mn}\,\KH(m/n)$ with $\KH$ given by the formula of Proposition~\ref{prop:kernel}: the
Gram matrix in the Lean statements is computed from the primes, and does not mention the zeros. \lean{BallWeil.form N
x} is $Q_N(x)$. The file \file{lib/check/Check.lean} prints the axioms of the theorems of Appendix~\ref{app:claims}, and
\texttt{scripts/check\_axioms.py} checks that they are \lean{propext}, \lean{Classical.choice} and \lean{Quot.sound} only.
The library's sources contain no \lean{sorry}, no declared axiom and no \lean{native\_decide}
(\texttt{scripts/check\_hygiene.py}). The interval certificate of Lemma~\ref{lem:minorant}(ii) and the computation
$\operatorname{lcm}(1, \dots, n) \le 3^n$ of Lemma~\ref{lem:mertens} are checked by kernel evaluation, \lean{decide
+kernel}, on exact arithmetic in natural numbers. The toolchain is Lean \texttt{v4.35.0-rc3} with Mathlib at commit
\texttt{0f64d30a}.

The statements that are not Lean theorems are: the remarks, among them the numerical values of
Section~\ref{sec:numerics} and the whole of Section~\ref{sec:analogues} except where a Lean theorem is named; the upper
bound for $d_N$ after Proposition~\ref{prop:growth}; the historical and comparative statements of
Sections~\ref{sec:intro} and~\ref{sec:related}; and the conjecture and questions of Section~\ref{sec:lambda}.

\section{The analogues}\label{sec:analogues}

This section collects tests of (A) and (B) on the natural analogues of the construction. Its statements are remarks:
they are proved here or computed, not formalised, except the density of real characters in Section~\ref{sec:chars},
which is a Lean theorem. The computations are in \texttt{scripts/ball\_tests/} of the repository, and
\file{docs/ball-conjecture-tests.md} gives the details.

\subsection{The first step}

\begin{remark}[The two forms at $k = 2$]
Every orbit Gram matrix of this kind has $G_{nn} = nG_{11}$, so $G_{22} = 2d_1$. Integer vectors see only $\mathrm{Re}\,G$,
since $Q(a) = a^T(\mathrm{Re}\,G)a$. With $g = \mathrm{Re}\,G_{12}/d_1$,
\[
Q(a\delta_1 + b\delta_2) = d_1\big((a + bg)^2 + b^2(2 - g^2)\big),
\]
so (B) holds at $N = 2$ exactly when $|g| \le 1$; otherwise $\delta_2 - \operatorname{sgn}(g)\delta_1$ has $Q = (3 -
2|g|)d_1 < d_1$. Since $d_2(\mathrm{Re}\,G) = (2 - g^2) d_1$, (B) at $N = 2$ is (A) at $k = 2$ for $\mathrm{Re}\,G$, and
for a real Gram matrix the two forms coincide on $w_1, w_2$. For $\zeta$, $g = -0.5306\ldots$
\end{remark}

\subsection{Curves over finite fields}

Let $C/\mathbb F_q$ have genus $g$ and $L$-polynomial $\prod_j (1 - \alpha_j T)$, $\alpha_j = \sqrt q\,e^{i\theta_j}$. Its
zeros are $\frac12 + i(\theta_j + 2\pi k)/\log q$, $k \in \mathbb Z$, and the same formula \eqref{eq:K} defines $\KH$
and $G$. Poisson summation gives $\KH(r) = \log q \sum_{m \in \mathbb Z} s_m q^{-m/2} e^{-|\log r - m \log q|/2}$, with
$s_m = \sum_j \alpha_j^m$ and $s_{-m} = q^{-m} s_m$, so every entry of $G$ is a rational multiple of $\log q$; the term
$m = 0$ alone is $2g \log q \min(m, n)$, the flat model.

\begin{remark}[The null vector is integral]
Let $m(z) = \sum_{k \le d} m_k z^{d - k}$ be the minimal polynomial of Frobenius, the squarefree part of $L^*(z) =
z^{2g} L(1/z) = \prod_j (z - \alpha_j)$: it is monic with integer coefficients, of degree $d \le 2g$. Then $\sum_k m_k
w_{q^{d - k}} = 0$. Indeed $q^{i\gamma_\rho} = e^{i\theta_j}$ depends only on the class $j$ of the zero, so $w_{q^a} =
\sum_j \alpha_j^a P_j w_1$, where $P_j$ projects onto the zeros of class $j$, and $\sum_k m_k w_{q^{d - k}} = \sum_j
m(\alpha_j) P_j w_1 = 0$. So $d_k = 0$ at $k = q^d$, and the integer vector with entry $m_k$ at $n = q^{d - k}$ has $Q
= 0$: both forms fail for every curve, at $N = q^d$. For the $148$ curves of genus $1$ to $3$ of the repository's sample,
$q \le 13$, this was checked in exact rational arithmetic; $q^d$ runs from $4$ to $13^6$.
\end{remark}

\begin{remark}[Below the null vector]
On $n \le \min(20, q^d - 1)$, (A) fails for $59$ of the $148$ curves, $43$ of them at $k = 2$, and (B) fails for exactly
those $43$, all at $N = 2$; on the first $20$ indices prime to $q$, where no null vector is forced, (A) fails for $53$
and (B) for $39$, again only at the first two indices. Failures were certified in exact rational arithmetic, passes by
a Fincke--Pohst search.
\end{remark}

\subsection{The Frobenius orbit}

The analogue that the conjecture's motivation starts from is the orbit of the graphs $\Gamma^k$ of the powers of
Frobenius on $C \times C$. Hallouin and Perret compute its Gram matrix, projected to the part orthogonal to the
fibres: $2gq^k$ on the diagonal and $q^{\min(i, k)} s_{|i - k|}$ off it, and read Weil's bound as a $2 \times 2$ Gram
determinant and Ihara's as a $3 \times 3$ one \cite[Lemma~3]{hallouin-perret}. In the variable $z$, the orbit is $1, z,
z^2, \dots$ with the form $T(P) = \sum_j |P(\alpha_j)|^2$.

\begin{remark}[The forms part]
(A) fails for every curve: the orbit vectors $z^0, \dots, z^d$ are dependent, so the pivot at $z^d$ is $0$; already
$d_2 \ge d_1$ means $|q + 1 - \#C(\mathbb F_q)| \le 2g\sqrt{q - 1}$, which maximal curves violate. But the integer span
of the orbit is the lattice $\mathbb Z[z]/(m)$, since the dependencies are integral, and $T$ is an integral form with
$T(1) = 2g$. If $m$ is irreducible, then $L^* = m^e$, and $\prod_j P(\alpha_j) = \operatorname{Res}(m, P)^e$ is a nonzero
integer for $P \notin (m)$, so by the arithmetic--geometric mean inequality $T(P) \ge 2g$: (B) holds. Equality at some
$P \ne \pm 1$ forces $P(\alpha)$ to be a root of unity other than $\pm 1$, as at $\alpha = 1 + i$ over $\mathbb F_2$,
where $P = z - 1$. In the sample, $m$ is irreducible for $132$ curves, and an exact search finds no failure of (B) for
the other $16$.
\end{remark}

\subsection{Dirichlet characters}\label{sec:chars}

For a primitive character $\chi$ mod $q$ of parity $a$, the same formula over the zeros of $L(s, \chi)$ defines
$\KH_\chi$ and $G$, and the explicit formula gives a closed form \cite[Chapter~19]{davenport}. With $K_1 = \KH_\chi(1)$
and $\ell = L'/L(1, \chi)$, $G_{12} = K_1 + 2c_a + \ell$, where $c_0 = -\frac14 \log\frac{27}{4}$ and $c_1 = \frac34
\log \frac34$.

\begin{remark}[The criterion at $k = 2$]
Since $K_1 = \log\frac q\pi + \psi_\Gamma(\frac{1 + a}2) + 2\operatorname{Re}\ell$, (B) holds at $N = 2$ if and only if
$\frac15 L_a \le K_1 \le L_a$, with $L_a = \log \frac q\pi + \psi_\Gamma(\frac{1 + a}{2}) - 4c_a$, for every $\chi$; for
real $\chi$ this is also (A) at $k = 2$. The upper end says that the low zeros carry more weight than their density
accounts for.
\end{remark}

\begin{remark}[Density]
Let $X$ be Lamzouri's random model for $L'/L(1, \chi_D)$ \cite{lamzouri}. The proportion of fundamental discriminants
$D > 1$ with $d_2 < d_1$ tends to $P(X > \frac12 \log\frac{27}{4}) = 0.0511\ldots$, and that of $D < 0$ to $P(X >
\frac32 \log \frac43) = 0.1609\ldots$ (\file{docs/d2-density.md}; Lean: \lean{D2Good.d2\_density\_pos\_closed},
\lean{D2Good.d2\_density\_neg\_closed}, whose proof uses the library's formalisation of a zero-free half-plane from an
unrefereed preprint, \file{formalisations/openai-half-plane.md}). By the remark at $k = 2$, these characters fail both
forms.
\end{remark}

\begin{remark}[The scans]
Of the $29468$ primitive characters with $q \le 400$, $5983$ fail (A) at $k = 2$ and $5443$ fail (B), all at $N = 2$,
on $n \le 12$; the other $540$ fail (A) only through $\operatorname{Im} \ell$. The four that first fail (A) at $k = 3$, two
mod $13$ and two mod $61$, pass (B). Of the $6086$ real primitive characters with $q \le 10^4$, (B) fails for exactly the
$637$ with $d_2 < d_1$. Failures were certified in ball arithmetic.
\end{remark}

\subsection{Li's basis}\label{sec:li}

\begin{remark}
Lagarias's Theorem~3.1 \cite{lagarias07} gives the Gram matrix of Weil's scalar product on Li's functions $G_n(s) = 1 -
(1 - 1/s)^n$: $\langle G_n, G_m\rangle_W = \lambda_n + \lambda_{-m} - \lambda_{n - m}$, with Li's coefficients
$\lambda_n$ \cite{li97}. With $F_n(s) = n^s/s$, the same pairing gives $G(m, n)$, and $G_1 = F_1 = 1/s$: the two are
Gram matrices of one scalar product on two orbits of one vector. Under RH, $z_\rho = 1 - 1/\rho$ lies on the unit
circle and the pivots of Lagarias's matrix are the squared norms of the monic orthogonal polynomials for $\sum_\rho
\delta_{z_\rho}/|\rho|^2$. They cannot increase, and they tend to $0$ by Szeg\H{o}'s theorem: numerically $0.0462$,
$7.42 \cdot 10^{-5}$, \dots, $5.89 \cdot 10^{-23}$ at $k = 8$ (\texttt{scripts/ball\_tests/prior\_art/li\_orbit.py}). In
the flat model both orbits have constant pivots; the zeros push the dilation orbit's pivots up and the shift orbit's
down.
\end{remark}

\section{Related work}\label{sec:related}

\emph{The positivity layer.} Weil's criterion \cite{weil52} in Bombieri's form \cite{bombieri00} is the starting point.
The kernel $\KH$ is Matsumoto and Suzuki's $H_1$ \cite{matsumoto-suzuki}, and their Theorem~1.3 is a continuous form of
Proposition~\ref{prop:positivity}; see also Suzuki's study of the screw function \cite{suzuki23}. Nicolas's bound
$|W(x)| \le \beta$ \cite{nicolas} is its $2 \times 2$ layer. The constant $d_1$ and the equivalence of RH with $\sum_\rho
1/|\rho|^2 = d_1$ are Burnol's \cite{burnol98} and Gun, Murty and Rath's \cite{gun-murty-rath}.

\emph{Weil's form on finite-dimensional spaces.} Yoshida \cite{yoshida92} proves non-degeneracy of Weil's form on
windows and, by Gram--Schmidt on high Fourier modes, explicit lower bounds for squared norms; Bombieri \cite{bombieri00,
bombieri03} studies its minimisation on windows. Connes and Consani \cite{connes-consani21, connes-consani23} study it on
$[\lambda^{-1}, \lambda]$, where its small eigenvalues are extremely small. None of these measures pivots against a first
pivot or considers integer lattices. The orbit vectors here are not compactly supported, and in Yoshida's setting only
the hyperplane $\sum_n x_n = 0$ of their span fits a window.

\emph{Distances and extremal problems.} In the Nyman--Beurling criterion, in Báez-Duarte's form \cite{baez-duarte03},
the squared distance from $1_{(0,1]}$ to the span of the first $N$ dilates tends to $0$; Burnol \cite{burnol02} bounds
it below by about $d_1/\log N$, and Bettin, Conrey and Farmer \cite{bcf} show, conditionally, that this is its size.
Here the squared distances of the $w_k$ from the span of the earlier ones are bounded below by $d_1$ instead. The
flat model's orthogonal polynomials are Lubinsky's \cite{lubinsky14}, and Dimitrov and Oliveira solve the dual extremal
problem, with the same constant \cite{dimitrov-oliveira}.

\emph{Analogues.} Hallouin and Perret's Gram matrix of the Frobenius orbit \cite{hallouin-perret} is the function-field
analogue of Section~\ref{sec:analogues}. Conrey and Li \cite{conrey-li} show that de Branges's positivity conditions,
stronger than RH, fail for $\zeta$; the margin here runs the other way, being equivalent to RH for $\zeta$ and false for
its analogues.

\emph{Li's criterion.} Li's coefficients \cite{li97} are values of Weil's form in another basis (Section~\ref{sec:li}).
Mazhouda \cite{mazhouda13} studies the monotonicity of Li's coefficients for Dirichlet $L$-functions; a monotonicity
statement equivalent to GRH would be a margin on the diagonal, not a margin for pivots.
% TODO: read Mazhouda (2013), Bombieri (2003) and Broughan's book before the novelty statements are final.

\section{The smallest eigenvalue, and questions}\label{sec:lambda}

The pivots of $G_N$ grow, but its smallest eigenvalue $\lambda_N = \min_{\|x\| = 1} Q_N(x)$ tends to $0$: the lattice
has a uniform gap at the integer points while the real quadratic form degenerates.

\begin{theorem}\label{thm:lambda}
Under RH, $\sqrt{\log N}/(220 e^8 N) \le \lambda_N$ for $N \ge 3$, and $\lambda_N \le 67(1 + \log N)/N$ for $N \ge 1$.
Without hypothesis, there is $C$ with $\lambda_N \le C(1 + \log N)/N$ for every $N \ge 1$.
\end{theorem}

The upper bounds take $x_n = (-1)^{n + 1}/n$, for which $\sum_n x_n n^\rho$ is a tail of the alternating zeta function
at $1 - \rho$, small at the zeros; the lower bound comes from a sampling theorem for Dirichlet polynomials at the zeros
up to height of order $N$ (\file{formalisations/ball-rate.md}).

\begin{conjecture}
$\lambda_N \sim c \log N/N$ with $c = 0.0429\ldots$
\end{conjecture}

\begin{remark}
Numerically $\lambda_N N = 0.0434 \log N + 0.032$ to within $2 \cdot 10^{-4}$ for $10^3 \le N \le 3 \cdot 10^4$
(\texttt{scripts/lambda\_large.py}), and a variational problem for the profile of the minimising vectors gives the constant
$0.0429$ (\file{scripts/packing.py}); an upper bound $\lambda_N \le (0.04294 + o(1)) \log N/N$ under RH is proved on paper
in \file{formalisations/ball-rate.md}, not in Lean. Neither the conjecture nor these values are Lean theorems.
\end{remark}

\begin{remark}[Questions]
(1) Is the second minimum of the lattice attained at $\pm(\delta_1 + \delta_2)$ for every $N$? Under RH it would follow
from $d_k \ge 1.939\,d_1$ for $k \ge 3$. (2) Is there an integral structure for $\zeta$ that explains (B), as integrality
explains it for the Frobenius orbit? (3) For Dirichlet $L$-functions, does a bounded number of pivots decide (A)?
\end{remark}

\section*{Acknowledgements}

This work was carried out with the assistance of AI systems, among them Claude, developed by Anthropic, and Grok Bot,
developed by xAI. They were used to develop the argument, to write the Lean formalisation and the numerical tests, and
to draft this paper. The author is responsible for the paper's content. The Lean verification is by the Lean kernel and
the repository's continuous integration, and not by the assistants.

\appendix

\section{The claims table}\label{app:claims}

The table names the Lean theorem behind each numbered result. A name's first component is its namespace, which is also
the name of its file, in \file{lib/core} except \lean{BallUpperEta}, which is in \file{lib/pnt}.
{\footnotesize
\begin{longtable}{@{}p{0.40\textwidth}>{\raggedright\arraybackslash}p{0.55\textwidth}@{}}
\toprule
Claim & Lean \\
\midrule
\endhead
Theorem~\ref{thm:main}, \ref{thm:main2}: RH $\iff$ (A) & \lean{PivotSharp.rh\_iff\_orbitLattice} \\
\quad (A) $\Rightarrow$ RH & \lean{BallWeil.rh\_of\_pivots} \\
\quad RH $\Rightarrow$ (A) & \lean{PivotSharp.orbitLattice\_of\_rh} \\
Theorem~\ref{thm:main}, \ref{thm:main2}: RH $\iff$ (B) & \lean{PivotSharp.rh\_iff\_shortest} \\
\quad RH $\Rightarrow$ (B) & \lean{PivotSharp.shortest\_of\_rh} \\
Theorem~\ref{thm:margin}: $Q_N(x) \ge c_0 x_N^2$ & \lean{PivotSharp.pivot\_ge\_rh} \\
\quad the integer vectors & \lean{PivotSharp.lattice\_gap} \\
Corollary~\ref{cor:lattice} & \lean{PivotSharp.shortest\_unique}, \lean{PivotSharp.lattice\_gap} \\
Proposition~\ref{prop:kernel} & \lean{BallHadamard.kerH\_one}, \lean{BallHadamard.kerH\_inv}, \lean{BallPsi.kerH\_explicit} \\
Proposition~\ref{prop:positivity} & \lean{BallWeil.rh\_iff\_form}, \lean{BallHadamard.rh\_iff\_realizationH}, \lean{BallWeil.form\_eq\_norm\_sq} \\
Proposition~\ref{prop:flat} & \lean{BallFlat.gram\_flat}, \lean{BallFlat.flat\_pivot} \\
Proposition~\ref{prop:growth} & \lean{BallPivot.pivot\_ge} \\
Lemma~\ref{lem:explicit} & \lean{WeilForm.weilExplicit\_zeta} \\
Lemma~\ref{lem:phi} & \lean{PivotPhi.integral\_phi\_cos}, \lean{PivotPhi.phi\_mul\_le}, \lean{PivotPhi.phiE\_strip}, \lean{PivotPhi.norm\_phiE\_I\_le} \\
Lemma~\ref{lem:prime} & \lean{PivotSharp.prime\_eq} \\
Lemma~\ref{lem:minorant}(i) & \lean{PivotMinorFT.integral\_mu\_cos\_eq\_zero}, \lean{PivotMinorFT.integral\_mu\_ge} \\
Lemma~\ref{lem:minorant}(ii) & \lean{PivotCert.mu\_le} \\
Lemma~\ref{lem:pole} & \lean{PivotSharp.pole\_ge} \\
Proposition~\ref{prop:form} & \lean{PivotSharp.form\_ge} \\
Lemma~\ref{lem:mertens} & \lean{PivotTop.mertens\_le\_uniform}, \lean{PivotTop.mertens\_le\_small} \\
Proposition~\ref{prop:top} & \lean{PivotTop.top\_bound} \\
Theorem~\ref{thm:lambda} & \lean{BallGap.rate\_gap}, \lean{BallUpperEta.rate\_upper\_rh}, \lean{BallUpperEta.rate\_upper} \\
\bottomrule
\end{longtable}
}

% The entries were taken from reports/Orbit lattice criterion prior art.md, which records how each was checked; Weil
% (1952) was not consulted, and Bombieri (2003), Lubinsky (2014) and Mazhouda (2013) are known only second-hand.
\begin{thebibliography}{99}
\raggedright
% The references of the orbit-lattice paper. The entries follow reports/Orbit lattice criterion prior art.md, which
% records how each was checked; the arXiv titles of Connes--Consani and Lamzouri were checked against arXiv on
% 13 November 2026, and Mazhouda's entry against Crossref. Not consulted: Weil (1952). Known only second-hand: Bombieri
% (2003), through its zbMATH review and Suzuki's citations, Lubinsky (2014), through Dimitrov and Oliveira, and
% Mazhouda (2013), through its abstract.
\bibitem{weil52} A. Weil, \emph{Sur les ``formules explicites'' de la th\'eorie des nombres premiers}, Comm. S\'em.
  Math. Univ. Lund, tome suppl\'ementaire (1952), 252--265.
\bibitem{matsumoto-suzuki} K. Matsumoto and M. Suzuki, \emph{M-functions and screw functions originating from
  Goldbach's problem and zeros of the Riemann zeta function}, J. Number Theory 280 (2026), 918--946;
  arXiv:2409.00888.
\bibitem{nicolas} J.-L. Nicolas, \emph{Small values of the Euler function and the Riemann hypothesis}, Acta Arith. 155
  (2012), 311--321; arXiv:1202.0729.
\bibitem{burnol98} J.-F. Burnol, \emph{The explicit formula in simple terms}, arXiv:math/9810169 (1998).
\bibitem{gun-murty-rath} S. Gun, M. R. Murty and P. Rath, \emph{Transcendental sums related to the zeros of zeta
  functions}, Mathematika 64 (2018), 875--897.
\bibitem{davenport} H. Davenport, \emph{Multiplicative Number Theory}, 2nd ed., Graduate Texts in Math. 74, Springer,
  1980.
\bibitem{lagarias07} J. C. Lagarias, \emph{Li coefficients for automorphic $L$-functions}, Ann. Inst. Fourier 57
  (2007), 1689--1740; arXiv:math/0404394.
\bibitem{li97} X.-J. Li, \emph{The positivity of a sequence of numbers and the Riemann hypothesis}, J. Number Theory 65
  (1997), 325--333.
\bibitem{yoshida92} H. Yoshida, \emph{On Hermitian forms attached to zeta functions}, in \emph{Zeta Functions in
  Geometry}, Adv. Stud. Pure Math. 21 (1992), 281--325.
\bibitem{bombieri00} E. Bombieri, \emph{Remarks on Weil's quadratic functional in the theory of prime numbers, I}, Rend.
  Mat. Acc. Lincei (9) 11 (2000), 183--233.
\bibitem{bombieri03} E. Bombieri, \emph{A variational approach to the explicit formula}, Comm. Pure Appl. Math. 56
  (2003), 1151--1164.
\bibitem{suzuki23} M. Suzuki, \emph{Aspects of the screw function corresponding to the Riemann zeta function}, J.
  London Math. Soc. (2) 108 (2023), 1448--1487; arXiv:2206.03682.
\bibitem{connes-consani21} A. Connes and C. Consani, \emph{Weil positivity and trace formula, the archimedean place},
  Selecta Math. (N.S.) 27 (2021), Paper 77; arXiv:2006.13771.
\bibitem{connes-consani23} A. Connes and C. Consani, \emph{Spectral triples and $\zeta$-cycles}, Enseign. Math. 69
  (2023), 93--148; arXiv:2106.01715.
\bibitem{baez-duarte03} L. B\'aez-Duarte, \emph{A strengthening of the Nyman--Beurling criterion for the Riemann
  hypothesis}, Rend. Mat. Acc. Lincei (9) 14 (2003), 5--11; arXiv:math/0202141.
\bibitem{burnol02} J.-F. Burnol, \emph{A lower bound in an approximation problem involving the zeros of the Riemann zeta
  function}, Adv. Math. 170 (2002), 56--70.
\bibitem{bcf} S. Bettin, J. B. Conrey and D. W. Farmer, \emph{An optimal choice of Dirichlet polynomials for the
  Nyman--Beurling criterion}, Proc. Steklov Inst. Math. 280 (2013), suppl. 2, 30--36; arXiv:1211.5191.
\bibitem{lubinsky14} D. S. Lubinsky, \emph{Orthogonal Dirichlet polynomials with arctangent density}, J. Approx.
  Theory 177 (2014), 43--56.
\bibitem{dimitrov-oliveira} D. K. Dimitrov and W. D. Oliveira, \emph{An extremal problem related to generalizations of
  the Nyman--Beurling and B\'aez-Duarte criteria}, arXiv:1608.07887.
\bibitem{hallouin-perret} E. Hallouin and M. Perret, \emph{A unified viewpoint for upper bounds for the number of points
  of curves over finite fields via Euclidean geometry and semi-definite symmetric Toeplitz matrices}, Trans. Amer. Math.
  Soc. 372 (2019), 5409--5451; arXiv:1409.2357.
\bibitem{conrey-li} J. B. Conrey and X.-J. Li, \emph{A note on some positivity conditions related to zeta and
  $L$-functions}, Internat. Math. Res. Notices 2000, no.~18; arXiv:math/9812166.
\bibitem{mazhouda13} K. Mazhouda, \emph{On the Li coefficients for the Dirichlet $L$-functions}, Monatsh. Math. 170
  (2013), 405--423.
\bibitem{lamzouri} Y. Lamzouri, \emph{The distribution of Euler--Kronecker constants of quadratic fields},
  arXiv:1410.1603.
\bibitem{lean4} L. de Moura and S. Ullrich, \emph{The Lean~4 theorem prover and programming language}, Automated
  Deduction -- CADE 28, Lecture Notes in Comput. Sci. 12699, Springer, 2021, 625--635.
\bibitem{mathlib} The mathlib Community, \emph{The Lean mathematical library}, Proceedings of the 9th ACM SIGPLAN
  International Conference on Certified Programs and Proofs (CPP 2020), 367--381.
\bibitem{repo} R. Comley, \emph{number-theory}, the repository of the paper and its Lean library,
  \url{https://github.com/dickie81/number-theory}. The paper describes the library at the commit that contains it.
\bibitem{repo-prior-art} \emph{Prior art for the orbit-lattice criterion}, the literature report
  \texttt{reports/Orbit lattice criterion prior art.md} of \cite{repo}, with its notes in
  \texttt{research\_notes/Orbit lattice criterion prior art/}.

\end{thebibliography}

\end{document}
